% Slide 10 – team, the ask to the jury and partners, closing. Team rows are set by \TeamRows in main-en.tex.
\begin{frame}{Team and what we need}
  \begin{columns}[T,onlytextwidth]
    \begin{column}{0.46\textwidth}
      \begin{block}{Team}
        \footnotesize
        \renewcommand{\arraystretch}{1.2}
        \begin{tabular}{@{}>{\bfseries}l l@{}}
          \TeamRows
        \end{tabular}
        \vspace{1mm}

        Prototype built during HackYeah; the code is in an open repository and runs with one command.
      \end{block}
    \end{column}
    \begin{column}{0.52\textwidth}
      \begin{block}{What we are looking for after the hackathon}
        \footnotesize
        \begin{itemize}
          \item \textbf{3~venues and 1~cultural institution} in Kraków for a 30-day pilot of the owner panel.
          \item \textbf{A community partner} (disability organisation) to verify data and test accessibility with users.
          \item \textbf{Contacts to a second city's open data} and to tourism organisations as a channel.
          \item \textbf{Pre-deployment funding} for 6~months (accelerator, accessibility grant) until Pro-plan
                revenue covers the team.
        \end{itemize}
      \end{block}
    \end{column}
  \end{columns}
  \vspace{2.5mm}
  \centering
  {\bfseries\color{accNavy}Not a static map: a living, verifiable picture of accessibility that pays for itself.\par}
  \vspace{0.8mm}
  \muted{Demo: \DemoUrl\ \textperiodcentered\ Code: \RepoUrl\ \textperiodcentered\ Video: \VideoUrl}

  \note[item]{Close with the sentence on this slide and invite the jury to the table with the phone: let them tap “Report a problem” themselves.}
  \note[item]{If asked about the funding amount: 6~months of work for 2--3 people plus infrastructure; base-scenario revenue appears from the second year.}
  \note[item]{Question about the team after the hackathon: an operator (company / foundation), moderation in the panel, funding from the Pro plan and white-label deployments.}
\end{frame}
